% ---------------------------------------------------------------------
%  Avant-propos + sommaire
% ---------------------------------------------------------------------
{\pltitre\fontsize{22}{26}\selectfont AVANT DE COMMENCER}\par
{\color{plOrange}\rule{\linewidth}{3pt}}\par

Vous avez un manuscrit prêt, ou une petite maison d’édition qui veut
publier sans immobiliser des milliers d’euros dans un stock ?
L’\textbf{impression à la demande} (IAD, ou \emph{print on demand}, POD)
est sans doute la solution. Ce guide fait le tour de ce qu’il faut savoir
en France : comment ça fonctionne, ce que dit la loi, combien ça coûte,
qui sont les acteurs, et comment s’y prendre étape par étape.
Et pour aller plus loin, des interviews d’auteurs et d’éditeurs vous
attendent sur la chaîne YouTube
\href{\plYoutubeURL}{\plYoutubeAffichage} (sélection en fin de guide).

\begin{astuce}[Pour qui ?]
\begin{itemize}
  \item \textbf{Auteurs indépendants} qui veulent un livre papier
        professionnel, vendu en ligne et, si possible, en librairie.
  \item \textbf{Petites maisons d’édition} qui veulent tester un titre,
        garder leur fonds disponible sans stock, ou rééditer des titres
        épuisés.
\end{itemize}
\end{astuce}

\begin{attention}[À lire]
Les tarifs, taux et conditions des plateformes changent souvent. Les
chiffres de ce guide sont des \textbf{ordres de grandeur} relevés en
2025-2026 : vérifiez-les toujours sur les sites officiels avant de vous
engager. Ce guide n’est pas un conseil juridique ou fiscal.
\end{attention}

\clearpage
\renewcommand{\contentsname}{Sommaire}
\setcounter{tocdepth}{2}
{\small\linespread{1}\selectfont\setlength{\parskip}{0pt}\tableofcontents}
